\makeabstract
    {
    Investigation of Intercellular Communication Between Mesenchymal Stromal Cells and Macrophages in a GUPy Hydrogel
    }
    {
    Raine Kamilova\textsuperscript{1}, Hana Yasue\textsuperscript{2,3}, Akihiro Nishiguchi\textsuperscript{2,3}
    }
    {
    \textsuperscript{1} Biochemistry \& Biophysics Program, Amherst College, Amherst, MA\\
\textsuperscript{2} Graduate School of Advanced Engineering, Tokyo University of Science, Tokyo, Japan\\
\textsuperscript{3} Research Center for Macromolecules and Biomaterials, National Institute for Materials Science, Tsukuba, Japan
    }
    {
    Mesenchymal stromal cells (MSCs) are the most widely used cell type in regenerative medicine and clinical trials due to their powerful secretome, and their tissue repairing and anti-inflammatory properties. Yet, these cells quickly succumb to apoptosis after transplantation, often breaking apart within 2 hours of injection. Although previous studies have shown the in-vivo efficacy of MSC transplantation via a GUPy hydrogel to boost cell survival, little was known about the role of hydrogels in regulating the MSC secretome and intercellular communications. Hence, this study aimed to identify how the hydrogel impacts MSC-macrophage cell communication and macrophage polarization.
    }
    {
    A UPy-modified gelatin (GUPy) hydrogel was formulated with sodium nicotinate (SN), which disrupted hydrogen bonds to allow for water solubility and subsequent gelation post-injection. The hydrogel{\textquoteright}s upper critical solution temperature (UCST) enabled high-temperature mixing and porous structure formation at body temperature. For indirect communication assays, macrophages were treated with lipopolysaccharide (LPS) and supernatant collected from 48-hour 2D or 3D (hydrogel-encapsulated) MSC cultures. Direct communication was assessed by co-culturing MSCs and macrophages for 48 hours, with LPS added at 24 hours. The experiments were evaluated using confocal laser microscopy, FIJI software, and ELISAs.
    }
    {
    TNF-\ensuremath{\alpha} ELISAs confirmed the significant anti-inflammatory effects of both 2D and 3D MSC cultures. Notably, 3D MSC supernatant (at 10\% and 25\% dilutions) significantly reduced pro-inflammatory (M1) macrophage polarization compared to 2D supernatant, implying that the hydrogel boosts MSCs' anti-inflammatory secretome. This was corroborated by increased IL-6 and HGF secretion from 3D-cultured MSCs. Furthermore, encapsulating macrophages alone in the hydrogel significantly reduced TNF-\ensuremath{\alpha}, suggesting inherent anti-inflammatory properties of the gel, likely due to SN. Finally, direct MSC-macrophage co-culture yielded significantly lower TNF-\ensuremath{\alpha} levels than indirect culture, highlighting that MSC-macrophage communication relies on more than just paracrine signaling.
    }
    {
    The GUPy hydrogel enhances the anti-inflammatory properties of encapsulated MSCs and exhibits intrinsic anti-inflammatory effects. Moreover, both direct intercellular contact and paracrine signaling collectively drive a decrease in the pro-inflammatory polarization of macrophages.
    }

    \newpage
    